% ============================================================
%  Chapitre 9 — Conclusion
% ============================================================
\chapter{Conclusion}

\section{Bilan}

Ce projet a permis d'implémenter et de comparer cinq formalismes de
représentation des connaissances et raisonnement, appliqués à deux domaines
réels : l'industrie musicale algérienne et le système judiciaire.

\medskip
\noindent Les objectifs initiaux ont été atteints :

\begin{enumerate}
  \item \textbf{Implémentation complète} des quatre formalismes pris en charge
        par le binôme : logique modale (TweetyProject), logique des défauts
        (TweetyProject RDL), logique classique (SAT4J + FOL) et logique de
        description (OWL API 4.5 + HermiT).

  \item \textbf{Démonstration de la non-monotonie} pour la logique des défauts
        (ajout de faits invalide des conclusions d'extension) et mise en
        évidence de l'hypothèse du monde ouvert pour la logique de description.

  \item \textbf{Interface graphique unifiée} permettant d'interagir avec chaque
        formalisme, d'observer les résultats du raisonneur en temps réel et de
        comparer les comportements.

  \item \textbf{Cohérence inter-formalismes} : les mêmes individus (Soolking,
        Ali, Karim…) sont représentés dans chaque formalisme, permettant une
        comparaison directe de leur expressivité sur des cas concrets.
\end{enumerate}

\section{Difficultés rencontrées}

\begin{description}
  \item[OWA en logique de description] L'hypothèse du monde ouvert d'OWL~2
        empêche de dériver l'absence d'une propriété à partir de son absence
        dans l'ABox. Pour les concepts définis par un complément (ex :
        ArtisteLibreDroits $\equiv$ Artiste $\dlinter$
        $\dlcomp{\dlsome{\text{signe\_chez}}{\text{Label}}}$),
        il a fallu recourir à des assertions explicites en ABox, analogues à
        une hypothèse de monde fermé locale.

  \item[Grounding dans Tweety RDL] Le calcul des extensions avec
        \code{SimpleDefaultReasoner} nécessite un grounding préalable.
        Avec plusieurs individus dans une même théorie, le \code{DefaultProcessTree}
        devient exponentiel. La solution a été de construire une théorie
        \emph{par individu}, ce qui rend le raisonnement tractable.

  \item[Compatibilité JavaFX] Le démarrage de JavaFX~17 avec Java~11 requiert
        une configuration spécifique du module-path Maven.
        Le plugin \code{javafx-maven-plugin} a été utilisé pour encapsuler
        cette configuration.

  \item[HermiT et les axiomes de disjonction] La vérification de disjonction
        entre deux classes via \code{isEntailed(OWLDisjointClassesAxiom)}
        requiert que la disjonction soit déductible structurellement — elle
        n'est pas déclarée implicitement, contrairement à l'intuition.
\end{description}

\section{Perspectives}

\begin{description}
  \item[Intégration SPARQL] Coupler la logique de description avec un point
        d'accès SPARQL (Apache Jena) pour permettre des requêtes expressive
        sur les ontologies OWL.

  \item[Raisonnement temporel] Étendre la logique modale avec des opérateurs
        temporels (CTL, LTL) pour modéliser l'évolution des carrières
        artistiques ou des procédures judiciaires dans le temps.

  \item[Interfaces Web] Remplacer l'interface JavaFX par une application Web
        (Spring Boot + React) pour rendre les formalismes accessibles depuis
        un navigateur.

  \item[Apprentissage automatique] Combiner la représentation symbolique
        des connaissances avec des techniques d'apprentissage automatique
        (neurosymbolique) pour extraire automatiquement des ontologies depuis
        des corpus textuels en langue arabe et en darija.

  \item[Logiques non-classiques supplémentaires] Implémenter la logique floue
        (pour le raisonnement incertain) ou la logique des croyances (pour
        la révision de croyances), qui complèteraient naturellement les
        formalismes déjà présents.
\end{description}

\bigskip
\noindent En définitive, ce projet illustre la richesse et la complémentarité
des formalismes de représentation des connaissances.
Aucun formalisme n'est universellement supérieur : le choix dépend du domaine,
de la nature des connaissances à représenter et des garanties de raisonnement
requises.
